% Generated by publication/build_sections.py; do not edit by hand.
\paragraph{Informal verdict.}
\textbf{A --- Complete and apparently correct}, with \textbf{high confidence (0.95)}.

\paragraph{Lean coverage.}
Q877.lean unconditionally proves the finite-dimensional iid integrable matrix strong law and an entrywise weak-law form, then the singular-mean obstruction; q877a gives a deterministic x != 0 annihilated by every A\_n on one probability-one event under almost-sure positive semidefiniteness; the explicit 3-by-3 Rademacher sequence is matrix-iid, integrable, entrywise bounded, centered, symmetric, rank 2 everywhere, has every partial sum singular, and rules out even an almost-sure kernel vector for A\_0; q877b derives the same nonzero simultaneous conclusion from a deterministic positive-definite H with H*A\_0 positive semidefinite almost surely, proving iid transfer, determinant/property-\(P\) preservation and cancellation by H internally. The source explicitly documents the printed omission of X != 0 and the p-to-N reindexing.

\subsection{Audited informal answer}
Q877
Restatement and a necessary correction
Let \((\Omega,\mathcal F,\mathbb P)\) be a probability space, let \(d\ge 1\), and let

\[
A_0,A_1,A_2,\ldots:\Omega\longrightarrow M_d(\mathbb R)
\]

be independent and identically distributed random matrices satisfying

\[
\mathbb E\|A_0\|<\infty
\]

for one, hence every, matrix norm. Put

\[
S_N:=A_0+\cdots+A_{N-1},\qquad N\ge1.
\]

Property \((P)\) is the existence of \(c\in(0,1)\) such that

\[
\mathbb P(\det S_N=0)\ge c
\qquad\text{for every }N\ge1.
\]

The questions are:


Prove that, when the \(A_n\) are symmetric positive semidefinite, property \((P)\) forces a common almost-sure kernel.


Find additional hypotheses under which the same conclusion remains valid more generally.


The PDF literally omits the condition \(X\ne0\). Without it, \(X=0\) makes both conclusions vacuous. I therefore solve the clearly intended statement with a deterministic nonzero

\[
X\in\mathbb R^d,\qquad
\mathbb P\bigl(A_nX=0\text{ for all }n\bigr)=1.
\]

Here ``symmetric positive'' is interpreted in the standard French sense: symmetric positive semidefinite. Luc Pommeret+1

Final results
The following conclusions hold.
Part \(a\): exact dichotomy
For iid positive semidefinite matrices, the first-moment assumption is in fact unnecessary. Exactly one of the following occurs:

\[
\boxed{\exists\,x\in\mathbb R^d\setminus\{0\}:
\mathbb P(A_0x=0)=1,}
\]

in which case

\[
\mathbb P(\det S_N=0)=1\qquad\forall N,
\]

or else

\[
\boxed{\mathbb P(\det S_N=0)\longrightarrow0.}
\]

Consequently, property \((P)\) is equivalent to the existence of a deterministic nonzero common null vector.
Part \(b\): failure in general, and an exact positive theorem
Symmetry, boundedness, centering, and constant rank do not suffice without positivity: there are bounded iid symmetric rank-\((d-1)\) matrices for which every \(S_N\) is singular but there is no common deterministic null vector.
On the other hand, the following natural additional hypothesis gives an exact affirmative answer:

If two independent copies commute almost surely,

\[
\mathbb P(A_0A_1=A_1A_0)=1,
\]

then, without any moment assumption,

\[
\limsup_{N\to\infty}\mathbb P(\det S_N=0)>0
\]

holds if and only if the matrices have a deterministic nonzero common null vector.

Thus pairwise commutation of the support is a complete sufficient structural hypothesis for part \(b\).

\begin{enumerate}
\item A general mean obstruction
\end{enumerate}
We first isolate the elementary probabilistic mechanism behind part \(a\).
Lemma 1
Let \(A_0,A_1,\ldots\) be iid integrable random matrices and let

\[
M:=\mathbb E A_0.
\]

If \(M\) is invertible, then

\[
\mathbb P(\det S_N=0)\longrightarrow0.
\]

Equivalently,

\[
\limsup_{N\to\infty}\mathbb P(\det S_N=0)>0
\quad\Longrightarrow\quad
\det(\mathbb E A_0)=0.
\]

Proof
We first recall the weak law of large numbers in the finite-dimensional matrix space:

\[
\frac{S_N}{N}\longrightarrow M
\qquad\text{in probability}.
\]

For completeness, this follows by truncation. Use the Frobenius norm and write

\[
A_k^{(R)}=A_k\mathbf 1_{\{\|A_k\|_F\le R\}},
\qquad
T_k^{(R)}=A_k-A_k^{(R)}.
\]

For fixed \(R\), independence and centering give

\[
\mathbb E\left\|
 \frac1N\sum_{k=0}^{N-1}
 \bigl(A_k^{(R)}-\mathbb EA_k^{(R)}\bigr)
\right\|_F^2
=
\frac1N
\mathbb E\left\|A_0^{(R)}-\mathbb EA_0^{(R)}\right\|_F^2,
\]

which tends to \(0\). Hence the bounded part converges in probability.
For the tail,

\[
\mathbb E\left\|
 \frac1N\sum_{k=0}^{N-1}
 \bigl(T_k^{(R)}-\mathbb ET_k^{(R)}\bigr)
\right\|_F
\le
2\mathbb E\|T_0^{(R)}\|_F.
\]

The right-hand side tends to \(0\) as \(R\to\infty\), by integrability. Markov's inequality completes the proof of the weak law.
Now suppose \(M\) is invertible. Since \(GL_d(\mathbb R)\) is open, there exists \(\delta>0\) such that

\[
\|B-M\|_F<\delta\quad\Longrightarrow\quad \det B\ne0.
\]

Because multiplication by \(N\ne0\) does not affect singularity,

\[
\begin{aligned}
\mathbb P(\det S_N=0)
&=\mathbb P\left(\det\frac{S_N}{N}=0\right)\\
&\le
\mathbb P\left(\left\|\frac{S_N}{N}-M\right\|_F\ge\delta\right)
\longrightarrow0.
\end{aligned}
\]

This proves the lemma. \(\square\)

\begin{enumerate}
\item Solution of part \(a\)
\end{enumerate}
Assume now that

\[
A_n=A_n^{\mathsf T}\succeq0
\qquad\text{almost surely}.
\]

By property \((P)\) and Lemma 1, the positive semidefinite matrix

\[
M=\mathbb EA_0
\]

is singular. Choose

\[
x\in\ker M\setminus\{0\}.
\]

Then

\[
0=x^{\mathsf T}Mx
  =\mathbb E\bigl[x^{\mathsf T}A_0x\bigr].
\]

The random variable \(x^{\mathsf T}A_0x\) is nonnegative. A nonnegative integrable random variable with expectation zero vanishes almost surely, so

\[
x^{\mathsf T}A_0x=0
\qquad\text{almost surely}.
\]

For a positive semidefinite matrix \(A\),

\[
x^{\mathsf T}Ax=\|A^{1/2}x\|^2,
\]

and therefore

\[
x^{\mathsf T}Ax=0\quad\Longrightarrow\quad Ax=0.
\]

Consequently,

\[
\mathbb P(A_0x=0)=1.
\]

The \(A_n\) have the same law, so

\[
\mathbb P(A_nx=0)=1
\qquad\forall n.
\]

Since the index set is countable,

\[
\mathbb P\bigl(A_nx=0\text{ for every }n\bigr)
=
\mathbb P\left(\bigcap_{n\ge0}\{A_nx=0\}\right)
=1.
\]

This proves part \(a\).

Stronger equivalence in the positive semidefinite case
For integrable iid positive semidefinite matrices, the following four statements are equivalent:

\[
\begin{array}{ll}
\text{(i)}&
\displaystyle\limsup_{N\to\infty}\mathbb P(\det S_N=0)>0,\\[1ex]
\text{(ii)}&
\mathbb EA_0\text{ is singular},\\[1ex]
\text{(iii)}&
\exists\,x\ne0:\ \mathbb P(A_0x=0)=1,\\[1ex]
\text{(iv)}&
\mathbb P(\det S_N=0)=1\quad\forall N.
\end{array}
\]

We proved \((i)\Rightarrow(ii)\Rightarrow(iii)\). If \(iii\) holds, then

\[
S_Nx=0
\]

almost surely, proving \(iv\); and \(iv\) trivially implies \(i\).
Thus, in the positive semidefinite setting, a uniform lower bound cannot take a genuinely intermediate value: it forces the singularity probability to be \(1\) for every partial sum.

The first moment is actually unnecessary in part \(a\)
For an arbitrary positive semidefinite random matrix \(A\), define

\[
B=\frac{A}{1+\|A\|}.
\]

Then \(B\) is bounded, positive semidefinite, and

\[
\ker B=\ker A.
\]

For positive semidefinite matrices \(C_1,\ldots,C_N\),

\[
\ker(C_1+\cdots+C_N)
=
\ker C_1\cap\cdots\cap\ker C_N.
\]

Indeed, if \((C_1+\cdots+C_N)x=0\), then

\[
0=x^{\mathsf T}(C_1+\cdots+C_N)x
 =\sum_{j=1}^N x^{\mathsf T}C_jx,
\]

and every summand is nonnegative, hence every \(C_jx=0\).
Consequently, if

\[
B_n=\frac{A_n}{1+\|A_n\|},
\]

then for every \(N\),

\[
\ker\sum_{n=0}^{N-1}A_n
=
\bigcap_{n=0}^{N-1}\ker A_n
=
\bigcap_{n=0}^{N-1}\ker B_n
=
\ker\sum_{n=0}^{N-1}B_n.
\]

Thus the singularity events are unchanged. The \(B_n\) are bounded, so the preceding proof applies to them. Hence the conclusion of part \(a\) holds without any moment assumption.
In particular,

\[
q_N:=\mathbb P(\det S_N=0)
\]

is nonincreasing in \(N\), and either \(q_N=1\) for every \(N\), or \(q_N\downarrow0\).

\begin{enumerate}
\item Why additional assumptions are indispensable in part \(b\)
\end{enumerate}
Let \((\xi_n,\eta_n)_{n\ge0}\) be iid pairs of independent Rademacher variables:

\[
\mathbb P(\xi_n=1)=\mathbb P(\xi_n=-1)=\frac12,
\qquad
\mathbb P(\eta_n=1)=\mathbb P(\eta_n=-1)=\frac12.
\]

Define

\[
A_n=
\begin{pmatrix}
0&\xi_n&\eta_n\\
\xi_n&0&0\\
\eta_n&0&0
\end{pmatrix}.
\]

These matrices are:


real symmetric;


bounded;


centered;


of rank \(2\) almost surely;


iid and possessing moments of every order.


For any \(x,y\in\mathbb R\), put

\[
M(x,y)=
\begin{pmatrix}
0&x&y\\
x&0&0\\
y&0&0
\end{pmatrix}.
\]

Every such matrix is singular. Moreover,

\[
S_N
=
M\left(\sum_{n=0}^{N-1}\xi_n,\,
       \sum_{n=0}^{N-1}\eta_n\right),
\]

so

\[
\mathbb P(\det S_N=0)=1
\qquad\forall N.
\]

Nevertheless, there is no nonzero deterministic common null vector. Indeed, if

\[
v=(v_1,v_2,v_3)^{\mathsf T}
\]

were annihilated almost surely, it would in particular be annihilated by the two support matrices \(M(1,1)\) and \(M(1,-1)\). The equations

\[
M(1,1)v=0,\qquad M(1,-1)v=0
\]

give

\[
v_1=0,\qquad v_2+v_3=0,\qquad v_2-v_3=0,
\]

hence \(v=0\).
Thus even the conjunction

\[
\text{symmetric}+\text{bounded}+\text{centered}+\text{fixed rank}
\]

does not suffice. The positivity hypothesis in part \(a\) is doing genuine work.

\begin{enumerate}
\item A precise affirmative answer to part \(b\): commuting matrices
\end{enumerate}
We prove the following exact theorem.
Theorem 2 --- commuting-support classification
Let \(A_0,A_1,\ldots\) be iid real \(d\times d\) random matrices. No moment assumption is required. Assume

\[
\mathbb P(A_0A_1=A_1A_0)=1.
\tag{C}
\]

Then the following are equivalent:


\(\displaystyle\limsup_{N\to\infty}\mathbb P(\det S_N=0)>0\);


there exists a deterministic \(x\in\mathbb R^d\setminus\{0\}\) such that

\[
\mathbb P(A_0x=0)=1;
\]



\(\mathbb P(\det S_N=0)=1\) for every \(N\).


In particular, under \(C\), property \((P)\) implies the desired conclusion.
The proof uses two elementary lemmas.

Lemma 2 --- atoms of sums of iid variables
Let \(Z_1,Z_2,\ldots\) be iid random variables with values in a finite-dimensional real vector space, or in \(\mathbb C\). Then either

\[
Z_1=0\quad\text{almost surely},
\]

or

\[
\mathbb P(Z_1+\cdots+Z_N=0)\longrightarrow0.
\]

Proof
Let \(\mu\) be the law of \(Z_1\), and let \(D\) be its set of atoms. The set \(D\) is countable. Put

\[
a=\mu(D).
\]

Case 1: \(a<1\)
Write

\[
\mu=a\mu_{\mathrm{at}}+(1-a)\mu_{\mathrm{na}},
\]

where \(\mu_{\mathrm{na}}\) has no atoms.
The convolution of an atomless probability measure with any probability measure is atomless, because

\[
(\mu_{\mathrm{na}}*\nu)(\{z\})
=
\int \mu_{\mathrm{na}}(\{z-y\})\,\nu(dy)
=0.
\]

Therefore an atom of \(\mu^{*N}\) can arise only from the term in which all \(N\) variables were sampled from the atomic component. Hence

\[
\sup_z\mu^{*N}(\{z\})\le a^N\longrightarrow0.
\]

Case 2: \(a=1\)
The law is purely atomic. If it is a point mass at \(z_0\ne0\), then

\[
Z_1+\cdots+Z_N=Nz_0\ne0.
\]

The only exceptional point mass is \(\delta_0\).
Otherwise choose two distinct atoms \(u\ne v\), of masses \(\alpha,\beta>0\). Set

\[
r=\alpha+\beta,\qquad q=\frac{\beta}{\alpha+\beta}\in(0,1).
\]

Let \(M_N\) be the number of indices \(i\le N\) for which \(Z_i\in\{u,v\}\). Then

\[
M_N\sim\operatorname{Bin}(N,r),
\qquad M_N\longrightarrow\infty
\quad\text{in probability}.
\]

Condition on the set of these indices and on all values outside \(\{u,v\}\). Given \(M_N=m\), the contribution of the selected variables is

\[
mu+K(v-u),
\qquad K\sim\operatorname{Bin}(m,q).
\]

Because \(v-u\ne0\), for any prescribed value of the total sum there is at most one possible integer \(K\). Consequently,

\[
\mathbb P(Z_1+\cdots+Z_N=0\mid M_N=m,\text{ other data})
\le
b_m(q),
\]

where

\[
b_m(q)=\max_{0\le k\le m}
\binom{m}{k}q^k(1-q)^{m-k}.
\]

Elementary Stirling bounds give

\[
b_m(q)=O_q(m^{-1/2}),
\]

so \(b_m(q)\to0\). Since \(M_N\to\infty\) in probability and \(0\le b_m(q)\le1\),

\[
\mathbb E[b_{M_N}(q)]\longrightarrow0.
\]

This proves the lemma. \(\square\)

Lemma 3 --- simultaneous triangularization
Every commuting family of complex \(d\times d\) matrices can be simultaneously upper triangularized.
Proof
First prove that a commuting family \(\mathcal A\) has a common eigenvector. Proceed by induction on the dimension.
If every member of \(\mathcal A\) is scalar, every nonzero vector works. Otherwise choose a nonscalar \(T\in\mathcal A\) and an eigenspace

\[
E_\lambda=\ker(T-\lambda I).
\]

This is a nonzero proper subspace. Since every \(B\in\mathcal A\) commutes with \(T\),

\[
B(E_\lambda)\subset E_\lambda.
\]

The restricted family on \(E_\lambda\) is commuting, so induction gives a common eigenvector.
After finding a common invariant line, pass to the quotient and iterate. This gives a common invariant flag and hence simultaneous upper triangularization. \(\square\)

Proof of Theorem 2
Let \(\mu\) be the common law and let

\[
K=\operatorname{supp}\mu.
\]

Step 1: the support is a commuting family
Assumption \(C\), together with independence, says

\[
(\mu\otimes\mu)\{(A,B):AB=BA\}=1.
\]

If two support points \(A,B\in K\) did not commute, continuity of the commutator would give neighborhoods \(U\ni A\), \(V\ni B\) such that no pair in \(U\times V\) commutes. Since \(A,B\) are support points,

\[
\mu(U)\mu(V)>0,
\]

contradicting \(C\). Thus every two matrices in \(K\) commute.
Let \(\mathcal A\) be the complex unital algebra generated by \(K\). It is commutative. By Lemma 3 there is \(P\in GL_d(\mathbb C)\) such that every

\[
P^{-1}AP,\qquad A\in\mathcal A,
\]

is upper triangular.
Write its diagonal entries as

\[
\chi_1(A),\ldots,\chi_d(A).
\]

Each \(\chi_j:\mathcal A\to\mathbb C\) is linear and multiplicative.
For every \(N\),

\[
P^{-1}S_NP
=
\sum_{n=0}^{N-1}P^{-1}A_nP
\]

is upper triangular, with diagonal

\[
\sum_{n=0}^{N-1}\chi_1(A_n),\ldots,
\sum_{n=0}^{N-1}\chi_d(A_n).
\]

Therefore

\[
\det S_N
=
\prod_{j=1}^d
\left(\sum_{n=0}^{N-1}\chi_j(A_n)\right).
\tag{1}
\]

Step 2: one joint character must vanish identically
By the union bound and \(1\),

\[
\mathbb P(\det S_N=0)
\le
\sum_{j=1}^d
\mathbb P\left(
\sum_{n=0}^{N-1}\chi_j(A_n)=0
\right).
\tag{2}
\]

For each \(j\), the variables \(\chi_j(A_n)\) are iid complex random variables. By Lemma 2, unless

\[
\chi_j(A_0)=0\quad\text{almost surely},
\]

the \(j\)-th probability on the right-hand side of \(2\) tends to zero.
If no \(\chi_j(A_0)\) were almost surely zero, the entire right-hand side of \(2\) would tend to zero, contradicting

\[
\limsup_N\mathbb P(\det S_N=0)>0.
\]

Hence there exists a joint character \(\chi\) among the \(\chi_j\) such that

\[
\chi(A)=0\qquad\text{for }\mu\text{-almost every }A.
\]

Since \(\chi\) is continuous, it vanishes on the support \(K\).
Step 3: a vanishing joint character yields a common kernel
Let \(\Psi\) be the finite set of distinct diagonal characters occurring in the simultaneous triangularization. Choose \(B\in\mathcal A\) such that

\[
\psi(B)\ne\chi(B)
\qquad\text{for every }\psi\in\Psi,\ \psi\ne\chi.
\]

This is possible because the forbidden set is a finite union of proper linear hyperplanes in \(\mathcal A\).
Let \(W\) be the generalized eigenspace of \(B\) corresponding to \(\chi(B)\):

\[
W=\ker(B-\chi(B)I)^d.
\]

It is nonzero because \(\chi\) occurs among the diagonal characters. Since every element of \(\mathcal A\) commutes with \(B\), the space \(W\) is invariant under \(\mathcal A\).
By the choice of \(B\), every diagonal character occurring on \(W\) is \(\chi\). Therefore, for every \(A\in K\), all eigenvalues of \(A|_W\) equal

\[
\chi(A)=0.
\]

Thus each \(A|_W\) is nilpotent. These restrictions form a commuting family of nilpotent matrices. By Lemma 3 they possess a common eigenvector \(w\in W\setminus\{0\}\). Since every member is nilpotent, every corresponding eigenvalue is zero:

\[
Aw=0\qquad\forall A\in K.
\]

The vector \(w\) may initially be complex. Write \(w=u+iv\), with \(u,v\in\mathbb R^d\). Since every \(A\in K\) is real,

\[
Au=0,\qquad Av=0.
\]

At least one of \(u,v\) is nonzero. Hence there exists a real nonzero vector \(x\) such that

\[
Ax=0\qquad\forall A\in K.
\]

In particular,

\[
\mathbb P(A_0x=0)=1.
\]

This proves \(1\Rightarrow2\). The implication \(2\Rightarrow3\) follows from

\[
S_Nx=0,
\]

and \(3\Rightarrow1\) is immediate. \(\square\)

Another broad sufficient hypothesis
Part \(a\) also immediately extends to matrices which are positive with respect to a common deterministic inner product.
Suppose there exists a deterministic positive definite matrix \(H\) such that

\[
HA_0
\]

is symmetric positive semidefinite almost surely. Then the matrices

\[
B_n=HA_n
\]

are iid positive semidefinite and

\[
\det\left(\sum_{n=0}^{N-1}B_n\right)
=
\det(H)\det\left(\sum_{n=0}^{N-1}A_n\right).
\]

Thus property \((P)\) is unchanged. Part \(a\) gives \(x\ne0\) with

\[
HA_nx=0
\]

almost surely for every \(n\), and the invertibility of \(H\) gives

\[
A_nx=0.
\]

So a common positive symmetrizer is another natural affirmative condition for part \(b\).

\begin{enumerate}
\item A general finite-variance classification of property \((P)\)
\end{enumerate}
There is a useful structural theorem explaining both the positive result and the counterexample.
Theorem 3 --- affine-hull dichotomy
Assume

\[
\mathbb E\|A_0\|^2<\infty.
\]

Let

\[
\mathcal H=\operatorname{aff}(\operatorname{supp}\mu)
\]

be the affine hull of the support of the law of \(A_0\). Then exactly one of the following occurs:


every matrix in \(\mathcal H\) is singular, in which case

\[
\mathbb P(\det S_N=0)=1
\qquad\forall N;
\]



\(\mathcal H\) contains an invertible matrix, in which case

\[
\mathbb P(\det S_N=0)\longrightarrow0.
\]



Consequently, under a finite second moment,

\[
\boxed{
(P)\iff
\operatorname{aff}(\operatorname{supp}\mu)
\subset\{A:\det A=0\}.
}
\]

Proof
Let

\[
m=\mathbb EA_0
\]

and write

\[
\mathcal H=m+V,
\]

where \(V\) is the direction space of the affine hull.
If every member of \(\mathcal H\) is singular, then

\[
\frac{S_N}{N}\in\mathcal H
\]

almost surely, since an affine subspace is closed under finite averages. Hence

\[
\det\frac{S_N}{N}=0,
\]

and therefore \(\det S_N=0\) almost surely.
Suppose now that the determinant does not vanish identically on \(\mathcal H\).
Define

\[
Z_N=\frac{S_N-Nm}{\sqrt N}.
\]

The multivariate central limit theorem, applied in the finite-dimensional vector space \(V\), gives

\[
Z_N\ \Longrightarrow\ G,
\]

where \(G\) is a centered Gaussian vector whose covariance is nondegenerate on \(V\). Indeed, if a nonzero linear functional on \(V\) had zero variance, it would vanish almost surely on \(A_0-m\), contradicting the minimality of the affine hull.
Expand the determinant along \(m+V\):

\[
\det(m+tz)=\sum_{r=0}^d t^rP_r(z),
\]

where each \(P_r\) is a homogeneous polynomial of degree \(r\) on \(V\). Since the determinant is not identically zero on \(m+V\), there is a least \(r_0\) for which

\[
P_{r_0}\not\equiv0.
\]

Because

\[
\frac{S_N}{N}=m+\frac{Z_N}{\sqrt N},
\]

we obtain

\[
N^{r_0/2}\det\left(\frac{S_N}{N}\right)
=
P_{r_0}(Z_N)
+
\sum_{r>r_0}
N^{-(r-r_0)/2}P_r(Z_N).
\]

The sequence \((Z_N)\) is tight, so every term in the sum tends to zero in probability. Hence

\[
N^{r_0/2}\det\left(\frac{S_N}{N}\right)
\Longrightarrow
P_{r_0}(G).
\]

A nonzero polynomial on a finite-dimensional real vector space has a Lebesgue-null zero set. This follows by induction on the dimension and Fubini's theorem: regarding the polynomial as a one-variable polynomial whose coefficients are polynomials in the other variables, outside the common zero set of the coefficients it has only finitely many roots. Since \(G\) has a density on \(V\),

\[
\mathbb P(P_{r_0}(G)=0)=0.
\]

Portmanteau's theorem now gives

\[
\begin{aligned}
\limsup_{N\to\infty}\mathbb P(\det S_N=0)
&=
\limsup_{N\to\infty}
\mathbb P\left(
N^{r_0/2}\det\left(\frac{S_N}{N}\right)=0
\right)\\
&\le
\mathbb P(P_{r_0}(G)=0)
=0.
\end{aligned}
\]

This proves the dichotomy. \(\square\)
The symmetric counterexample above has an affine hull equal to the two-dimensional linear space

\[
\left\{
\begin{pmatrix}
0&x&y\\
x&0&0\\
y&0&0
\end{pmatrix}:x,y\in\mathbb R
\right\},
\]

which consists entirely of singular matrices but has trivial common kernel. This is precisely the obstruction in the general problem.

\begin{enumerate}
\item Boundary cases and the role of the hypotheses
\end{enumerate}
For \(d=1\), part \(a\) reduces to a simple scalar statement. If \(A_n\ge0\), then

\[
A_0+\cdots+A_{N-1}=0
\]

if and only if every \(A_n=0\). Thus

\[
\mathbb P(S_N=0)=\mathbb P(A_0=0)^N.
\]

A positive uniform lower bound is possible only when \(A_0=0\) almost surely.
If the \(A_n\) are positive definite almost surely, every finite sum is positive definite, so property \((P)\) is impossible.
The upper restriction \(c<1\) is immaterial: every proof only uses \(c>0\). Likewise, ``for every \(N\)'' can be weakened to

\[
\limsup_{N\to\infty}\mathbb P(\det S_N=0)>0.
\]

Independence, or at least an ergodic law of large numbers, is genuinely relevant. For example, choose one random variable \(J\in\{1,2\}\), uniformly, and set for every \(n\)

\[
A_n=
\begin{cases}
\operatorname{diag}(0,1),&J=1,\\
\operatorname{diag}(1,0),&J=2.
\end{cases}
\]

The sequence is identically distributed, bounded, and positive semidefinite, and every partial sum is singular. But

\[
\ker\operatorname{diag}(0,1)
\cap
\ker\operatorname{diag}(1,0)
=\{0\},
\]

so there is no deterministic common nonzero null vector. This sequence is perfectly dependent.

Conclusion
Under the intended nonzero-vector interpretation, part \(a\) admits a complete and stronger answer:

\[
\boxed{
\text{For iid positive semidefinite matrices, }
(P)
\iff
\exists\,x\ne0:\ A_nx=0\text{ a.s. for all }n.
}
\]

No moment assumption is actually required, and if the common kernel does not exist then

\[
\mathbb P(\det S_N=0)\to0.
\]

For part \(b\), the conclusion is false even for bounded centered symmetric rank-\((d-1)\) matrices. A precise natural affirmative answer is obtained by requiring the support to be commuting:

\[
\boxed{
\mathbb P(A_0A_1=A_1A_0)=1
\quad\Longrightarrow\quad
(P)
\iff
\exists\,x\ne0:\ A_nx=0\text{ a.s. for all }n.
}
\]

Under a finite second moment, property \((P)\) has the additional exact algebraic characterization that the affine hull of the law's support consists entirely of singular matrices.
